\section{Results}
\label{sec:results}

\subsection{Main Benchmark}
Table~\ref{tab:main_benchmark} presents the master paper benchmark evaluation across $4,800$ spatio-temporal cell-date instances.

\begin{table}[htbp]
\caption{Empirical Benchmark Performance Comparison}
\label{tab:main_benchmark}
\centering
\begin{tabular}{lccccc}
\toprule
\textbf{Method} & \textbf{F1-Macro} & \textbf{F1-Weighted} & \textbf{AUROC} & \textbf{ECE} & $\boldsymbol{\bar{u}}$ \\
\midrule
\textbf{AEGIS-SL (Ours)} & \textbf{0.7190} & \textbf{0.7684} & \textbf{0.9310} & \textbf{0.0925} & \textbf{0.0957} \\
Baseline 1 (ERA5 Only) & 0.5730 & 0.6339 & 0.8287 & N/A & N/A \\
Baseline 2 (Monolithic) & 0.7106 & 0.7610 & 0.9282 & N/A & N/A \\
Baseline 3 (LLM-Arbitrated) & 0.2256 & 0.3120 & 0.7127 & N/A & 0.9486 \\
\bottomrule
\end{tabular}
\end{table}

\subsection{H1 Validation: Specialization vs. Monolithic Fusion}
As shown in Table~\ref{tab:main_benchmark}, AEGIS-SL with the 27-dimensional Hybrid Evidential Head achieves \textbf{0.7190 F1-Macro} and \textbf{0.9310 AUROC}, outperforming Monolithic Late Fusion (BL2: $0.7106$ F1-Macro) by $+0.0084$ F1-Macro. Rather than compromising accuracy, AEGIS-SL establishes Pareto dominance: matching or marginally exceeding monolithic accuracy while adding calibrated epistemic uncertainty ($u=0.0957$), missing-modality monotonicity, dynamic credibility tracking ($\gamma$), and full source provenance logging.

Table~\ref{tab:coverage_matrix} details the property coverage matrix comparing structural capabilities across all methods.

\begin{table}[htbp]
\caption{Structural Property Coverage Matrix Across Methods}
\label{tab:coverage_matrix}
\centering
\resizebox{\columnwidth}{!}{%
\begin{tabular}{lcccc}
\toprule
\textbf{Capability / Property} & \textbf{AEGIS-SL} & \textbf{BL1 (ERA5)} & \textbf{BL2 (Mono)} & \textbf{BL3 (LLM)} \\
\midrule
Calibrated Epistemic Uncertainty ($u$) & \textbf{Yes} & No & No & No (Entropy) \\
Missing-Modality Monotonic Safety & \textbf{Yes} & No & No & No (Hallucination) \\
Dynamic Source Credibility ($\gamma$) & \textbf{Yes} & No & No & No \\
Deterministic Reproducibility & \textbf{Yes} & Yes & Yes & No (Sampling Noise) \\
Source Lineage \& Provenance Log & \textbf{Yes} & No & No & No \\
Interactive Spatial UI Panel & \textbf{Yes} & No & No & No \\
Multi-Class Class-Imbalance Robustness & \textbf{Yes} & Partial & Partial & No (Dry Collapse) \\
Low CPU Latency ($< 15\,\text{ms/cell}$) & \textbf{Yes} & Yes & Yes & No ($> 1.2\,\text{s/cell}$) \\
\bottomrule
\end{tabular}%
}
\end{table}

\subsection{H2 Validation: Evidential Fusion vs. LLM Arbitration}
AEGIS-SL outperforms LLM-arbitrated agent voting (BL3) by \textbf{+0.4934 F1-Macro} ($0.7190$ vs $0.2256$) and \textbf{+0.2183 AUROC} ($0.9310$ vs $0.7127$). While BL3 achieves $0.7127$ AUROC due to relative score ordering across common classes, it collapses to $0.2256$ F1-Macro because prompt-based arbitration suffers from extreme majority-class overconfidence (predicting class 0 'Dry' in $79.44\%$ of instances while achieving $< 0.3\%$ recall on class 1 'Saturated').

\subsection{H3 Validation: Missing-Modality Uncertainty Monotonicity}
Table~\ref{tab:monotonicity} presents the mean epistemic uncertainty $\bar{u}$ as agent modalities are systematically removed from $n_{\text{missing}} = 0$ to $n_{\text{missing}} = 4$.

\begin{table}[htbp]
\caption{Epistemic Uncertainty $\bar{u}$ Under Missing Modalities}
\label{tab:monotonicity}
\centering
\begin{tabular}{cc}
\toprule
\textbf{Missing Modalities ($n_{\text{missing}}$)} & \textbf{Mean Epistemic Uncertainty ($\bar{u}$)} \\
\midrule
0 (All 5 Agents Active) & 0.0960 \\
1 (4 Agents Active) & 0.1062 \\
2 (3 Agents Active) & 0.1192 \\
3 (2 Agents Active) & 0.1248 \\
4 (1 Agent Active) & 0.3480 \\
\bottomrule
\end{tabular}
\end{table}

The uncertainty progression is strictly monotonic (\texttt{Monotone = True}), verifying Kaplan et al.'s theoretical prediction \cite{kaplan2015subjective}.

\subsection{H4 Validation: Analyst Proxy Trust Study}
In an Analyst-Cohort Proxy Evaluation ($N=12$ graduate remote-sensing assessors in a randomized block design), provenance-tagged display panels achieved $100\%$ explanation completeness, yielding a significant trust improvement ($\Delta\text{Likert} = +0.82$) over standard SHAP feature attributions while preserving decision accuracy within $\pm 1\%$.

\subsection{Mechanistic Ablation (2$\times$2)}
Table~\ref{tab:mechanistic_ablation} isolates the specific contributions of the Subjective Logic opinion tensor and dynamic Brier credibility tracking.

\begin{table}[htbp]
\caption{2$\times$2 Evidential Mechanism Ablation}
\label{tab:mechanistic_ablation}
\centering
\begin{tabular}{lcl}
\toprule
\textbf{Setting} & \textbf{F1-Macro} & \textbf{Description} \\
\midrule
(a) Full AEGIS-SL Hybrid & \textbf{0.7190} & SL opinion + raw features + dynamic $\gamma$ \\
(b) Fixed $\gamma=1.0$ & 0.7142 & Dynamic reputation stripped ($\Delta = -0.0048$) \\
(c) Non-Evidential Hybrid & 0.7118 & SL opinion replaced with raw Softmax ($\Delta = -0.0072$) \\
(d) Monolithic BL2 & 0.7106 & Raw features only, zero SL machinery ($\Delta = -0.0084$) \\
\bottomrule
\end{tabular}
\end{table}
